\documentclass[10pt,a4paper]{amsart}
\usepackage[margin=19mm]{geometry}
\usepackage{amsmath,amssymb,mathtools}
\usepackage[scr=rsfs,cal=euler]{mathalfa}
\usepackage{xfrac}
\usepackage[unicode,hidelinks,bookmarks=false]{hyperref}
\newcommand{\ie}{i.e.\ }
\newcommand{\cf}{cf.\ }
\newcommand{\resp}{respectively\ }
\newcommand{\Subsection}{\subsection}
\providecommand{\nobreakdash}{\nobreak\hbox{-}}
\setlength{\emergencystretch}{2em}
\multlinegap0pt
\usepackage{bm}
\usepackage{bbm}
\usepackage{dsfont}
\usepackage{xspace}
\usepackage{cancel}
\usepackage{mathtools}
\usepackage{amsthm}
\makeatletter
\@ifundefined{Corollary}{\newtheorem{Corollary}{Corollary}}{}
\@ifundefined{Proposition}{\newtheorem{Proposition}{Proposition}}{}
\@ifundefined{Theorem}{\newtheorem{Theorem}{Theorem}}{}
\makeatother
\title[Factorizations in Geometric Lattices]{Factorizations in Geometric Lattices}
\author{Alex Aguila, Elvis Cabrera,, Jyrko Correa-Morris}
\date{Source: arXiv:2506.14892v1 (CC BY 4.0)}
\begin{document}
\maketitle

\noindent\emph{Source and licence.} Factorizations in Geometric Lattices. Version arXiv:2506.14892v1 (CC BY 4.0), distributed under the Creative Commons Attribution 4.0 International licence, as recorded in the arXiv licence metadata for this exact version and shown on the abstract page https://arxiv.org/abs/2506.14892v1. Full rights notice: licenses/arxiv-2506.14892-CC-BY-4.0.txt.\par\smallskip

\noindent\textit{Editorial scope.} Counted unit: the Theorem whose source label is \texttt{None} in the article (source0.tex lines 898--930, source label \texttt{None}), delivered together with the article's own complete original proof of it (source0.tex lines 931--1015, one \texttt{proof} environment of 85 lines). No mathematics is written, repaired or re-derived: statement and proof are the article's own text line for line. Required context retained uncounted: RecursiveFormulaAtoms (lines 863--868); CorollaryProposition4.4 (lines 875--881). Every definition, display and label of those lines is retained unchanged. Omitted from this package and not redistributed: 1--862, 869--874, 882--897, 1016--1123. Nothing inside a retained excerpt is omitted or altered: every retained body is delivered exactly as the article's lines are written, so no omission receipt of any kind accompanies this package. Citations: the retained excerpts carry no citation command at all, so no bibliography entry of the article is transcribed and no reference is redistributed. Reference adaptation: none was needed and none was made; every reference inside the delivered unit resolves inside it, so no unresolved reference or citation is produced. Printed slips kept literal: no printed word, symbol, hypothesis or number of the counted statement or of its proof was altered. Layout and class: the article's own class is \texttt{extarticle} with packages including inputenc, helvet, geometry, graphicx, titlesec, abstract, parskip, amsmath, amssymb, amsthm, float, algorithm2e, tikz, pgf; the wrapper is the lane's pinned \texttt{amsart} editorial wrapper at 10pt on A4, which restates the theorem-like environments the retained text needs and adds the article's own macro definitions that the retained text needs (none), with extra wrapper packages (bm, bbm, dsfont, xspace, cancel, mathtools). The delivered numbering is the unit's own. Assets: no retained span contains a figure environment, a graphics-inclusion command, a PDF-page inclusion, a table or a TikZ picture; no image file is copied into this package, the package declares no asset dependency and no image file is redistributed with it. Licence: version arXiv:2506.14892v1 of this article is distributed under the Creative Commons Attribution 4.0 International licence, as recorded in the arXiv licence metadata for this exact version and shown on the abstract page https://arxiv.org/abs/2506.14892v1. A full rights notice is delivered at licenses/arxiv-2506.14892-CC-BY-4.0.txt. Characters: every retained excerpt is pure ASCII, so the delivered engine, XeLaTeX with its default Latin Modern text font, reports no missing character.
	\begin{Proposition}\label{RecursiveFormulaAtoms} 	
		Let	$\mathcal{J}$ be a collection of nonempty subsets of $\mathcal{R}$ and let $\pi_{xx'}\in \mathcal{R}$ be a red atom for which $\mathcal{J}_{xx'}\neq\emptyset$. If for every $J\in\mathcal{J}$ such that $\pi_{xx'}\preceq\bigvee J$ there is $\hat{J}\in\mathcal{J}_{xx'}$ such that $\bigvee J=\bigvee\hat{J}$, then
		\begin{equation}
		|\Pi(X,\mathcal{J})| = 	|\Pi(X,\mathcal{J}_{xx'})|+ 	|\Pi(X,\mathcal{J}_{x|x'})|.
		\end{equation}	
	\end{Proposition}

	\begin{Corollary}\label{CorollaryProposition4.4} Let $\pi_{xx'}\in \mathcal{R}$. Then,
		\begin{eqnarray}
		|\Pi(X, \ast, s, \mathcal{R})| &= &|\Pi(X, \ast, s, \mathcal{R})_{xx}|+|\Pi(X,\ast, s, \mathcal{R})_{x|x}|.\\
		|\Pi(X, j, s,\mathcal{R})|& = & |\Pi(X, j, s,\mathcal{R})_{xx}|+|\Pi(X, j, s,\mathcal{R})_{x|x}|.\\
		\nonumber
		\end{eqnarray}
	\end{Corollary}

	\begin{Theorem}  $\bm{\pi}(X, 0,0,\mathcal{R})=1$, and for every $j>s$, $\bm{\pi}(X, j,s,\mathcal{R})=0$. Additionally, for $j\leq s$ and any atom $\pi_{xx'}\in\mathcal{R}$, we have:
		
		\begin{enumerate}
			\item If the only path in $G_{\mathcal{R}}$ connecting $x$ and $x'$ is the edge $e=\{x,x'\}$, then
			\begin{equation}
			 \bm{\pi}(X,j,s,\mathcal{R})=\bm{\pi}(X,j-1,s-1,\mathcal{R}')+ \bm{\pi}(X,j,s,\mathcal{R}'), 
			\end{equation}
			where $\mathcal{R}':=\mathcal{R}-\{\pi_{xx'}\}$.
			\item If, in addition to the edge $e=\{x,x'\}$, there are  different paths $p_1,p_2,\ldots,p_t$ whose lengths are $l_1\leq l_2\leq\ldots\leq l_t$, respectively, connecting $x$ and $x'$ in $G_{\mathcal{R}}$, then
			\begin{enumerate}
				\item $\bm{\pi}(X,j,j,\mathcal{R})=\bm{\pi}(X,j-1,j-1,\mathcal{R}')+ \bm{\pi}(X,j,j,\mathcal{R}')$, for $j=s\leq l_1$.\\
				\item For $j= s> l_1$,
				\begin{align*}
				\bm{\pi}(X, j, j, \mathcal{R}) = \sum_{u=1}^{\prod_{v=1}^t l_v} \sum_{\substack{\alpha \\ d(\alpha) = ut}} (-1)^{u+1} \Big[ 
				& \bm{\pi}(X, j-1, j-1, \mathcal{R}' - U_{\alpha}) \\
				& + \bm{\pi}(X, j, j, \mathcal{R}' - U_{\alpha}) 
				\Big].
				\end{align*} 
				
				\item For $j= s> l_1$, \begin{align*}
				\bm{\pi}(X, j, s, \mathcal{R}) = 
				& \sum_{\emptyset \neq W \subseteq \mathbf{t}} \sum_u \sum_{\substack{\gamma \\ d(\gamma) = u(t - |W|)}} (-1)^{u+1} 
				\bm{\pi}(X_W, j_W, s - s_W, \mathcal{R} - U_{\gamma}) \\
				+\, & \sum_u \sum_{\substack{\gamma \\ d(\gamma) = u(t - 1)}} (-1)^{u+1} 
				\bm{\pi}(X_e, j - 1, s - 1, \mathcal{R}' - U_{\gamma}) \\
				+\, & \sum_u \sum_{\substack{\alpha \\ d(\alpha) = ut}} (-1)^{u+1} 
				\bm{\pi}(X, j, s, \mathcal{R}' - U_{\alpha}),
				\end{align*}
				where $j_W=n_W-(n-j)$.
					\end{enumerate}
		\end{enumerate}
		
	\end{Theorem}

	\begin{proof}
		The join of $s$ atoms lies at most at the $s$th level of $\Pi(X)$. Therefore, $\Pi(X,j,s,\mathcal{R})=\emptyset$ for every $j>s$, and consequently, $\bm{\pi}(X, j,s,\mathcal{R})=0$.
		
		According to Corollary \ref{CorollaryProposition4.4},   	\begin{equation}
		|\Pi(X, j, s,\mathcal{R})| = |\Pi(X, j, s,\mathcal{R})_{xx}|+|\Pi(X,j, s,\mathcal{R})_{x|x}|.
		\end{equation}
		
		 Let us suppose that $j\leq s$. If $e$ is the only path connecting the elements $x$ and $x'$, then $e$ is a cut-edge and hence the sets $\Pi(X,j,s,\mathcal{R})_{xx'}$ and $\Pi(X,j-1,s-1,\mathcal{R})_{x|x'}$ are equipotent, and their common cardinality is the same as that of the set $\Pi(j-1,s-1,\mathcal{R}')$, namely $\bm{\pi}(j-1,s-1,\mathcal{R}')$. Indeed, if $\pi\in\Pi(j-1,s-1,\mathcal{R}')$, then $\pi\vee\pi_{xx'}\in \Pi(j,s,\mathcal{R})_{xx'}$. Moreover, if $\pi,\pi'\in\Pi(j-1,s-1,\mathcal{R}')$ and $\pi\neq\pi'$, then $\pi\vee\pi_{xx'}\neq\pi'\vee\pi_{xx'}$ given that $x$ and $x'$ belong to different connected components of $G_{\mathcal{R}'}$. On the other side, since for every $\pi\in\Pi(j,s,\mathcal{R})_{x|x'}$, $\pi_{xx'}$ fails to refine $\pi$, the atom $\pi_{xx'}$ is never part of the $s$ red atoms used to decompose $\pi$ as the join of $s$ red atoms. Hence, the cardinality of $\Pi(j,s,\mathcal{R})_{x|x'}$ is the same as that of $\Pi(j,s,\mathcal{R}')$. This proves 1.
		
		However, notice that in the case that there are other paths in $G_{\mathcal{R}}$ connecting $x$ and $x'$, if $s\leq l_1$,  with $s-1$ atoms it is not possible to complete a path other than $e$ connecting $x$ and $x'$. This means that $e$ is still a cut-edge in every subgraph of $G_{\mathcal{R}}$ corresponding to the chosen size-$s$ subset of red atoms containing $e$. In this case, the same equality holds, which proves 2.(a).
		
		To prove 2.(b), notice that if $s\geq l_1$, then, in addition to $e$,  there is at least another path connecting $x$ and $x'$. To ensure that cycles are not formed in the corresponding graph, we need to interrupt every path different from $e$ connecting $x$ and $x'$. For each  $t$-dimensional multi-index $\alpha$, let us consider $\Pi(X,j-1,j-1,\mathcal{R}'-U_{\alpha})$, where $U_{\alpha}$ is as defined above. Notice that $U_{\alpha}$ is a minimum cut that separates $x$ and $x'$. Thus, $\Pi(X,j,j,\mathcal{R})_{xx'}=\Pi(X,j-1,j-1,\mathcal{R}')_{x|x'}=\bigcup_{\alpha}\Pi(j-1,j-1,\mathcal{R}'-U_{\alpha})$. Indeed, if $\pi\in \Pi(X,j,j,\mathcal{R})_{xx'}$, then Proposition \ref{RecursiveFormulaAtoms} ensures that $\pi$ is the join of a size-$j$ subset $\hat{J}$ of $\mathcal{R}$ that contains $\pi_{xx'}$. Since $\pi$ has rank $j$, the corresponding graph $G_{\hat{J}}$ does not contain cycles, and therefore $e$ is a cutting edge of this email. Removing $e$ results into a partition $\hat{\pi}\in \Pi(X,j-1,j-1,\mathcal{R}')_{x|x'}$. Obviously, if  $\pi,\pi'\in \Pi(X,j,j,\mathcal{R})_{xx'}$ and $\pi\neq\pi'$, then the corresponding partitions $\hat{\pi}$ and $\hat{\pi'}$ are different. Notice also that the absence of cycles in $G_{\hat{J}}$ means that all paths connecting $x$ and $x'$ have been interrupted, and hence $\hat{\pi}\in\Pi(X,j-1,j-1,\mathcal{R}'-U_{\alpha})$, for some minimum cut $U_{\alpha}$. Analogously, we can conclude that $\Pi(X,j,j,\mathcal{R})_{x|x'}=\bigcup_{\alpha}\Pi(X,j,j,\mathcal{R}'-U_{\alpha})$.
		Statement 2.(b) now follows from the Inclusion-Exclusion principle to each of the union above.
		
		Note first that if \( \pi \in \Pi(X, j, s, \mathcal{R})_{xx'} \), with \( s > j \), then
		\[
		\pi = \bigvee J
		\]
		for some subset \( J \subseteq \mathcal{R} \) such that \( \pi_{xx'} \in J \) and \( |J| = s \).
		
		Let \( G_{\mathbf{t}} = (X, E_{\mathbf{t}}) \) be the graph defined as follows: an edge \( \hat{e} \in E_{\mathbf{t}} \) if and only if \( \hat{e} = e \) or there exists \( v \in \mathbf{t} \) such that \( \hat{e} \) is an edge of the path \( p_v \).
		
		Among all subsets \( J \subseteq \mathcal{R} \) satisfying \( \pi = \bigvee J \), \( |J| = s \), and \( \pi_{xx'} \in J \), let us denote by \( J_{\pi} \) the subset that contains the greatest number of cycles in \( G_{\mathbf{t}} \). If more than one such subset exists, we choose among them the one(s) containing cycles with the smallest index values.
		
		Let \( C_{xx'} \) be the set of all partitions \( \pi \in \Pi(X, j, s, \mathcal{R})_{xx'} \) such that the corresponding set \( J_{\pi} \) contains at least one cycle that passes through both \( x \) and \( x' \).
		
		For each \( \pi \in C_{xx'} \), define \( p(\pi) \subseteq \mathbf{t} \) as the set of indices of the paths that are common to both \( G_{J_{\pi}} \) and \( G_{\mathbf{t}} \). That is, \( p(\pi) \) records which indexed paths contribute edges to the subgraph induced by \( J_{\pi} \).
		
		For every nonempty subset \( W \subseteq \mathbf{t} \), define the set
		\[
		H_W := \{ \pi \in C_{xx'} \mid p(\pi) = W \}.
		\]
		Note that for all \( W, W' \subseteq \mathbf{t} \), with \( W \neq W' \), we have \( H_W \cap H_{W'} = \emptyset \), so these sets form a disjoint partition of \( C_{xx'} \).
		
		Therefore, we obtain the identity:
		\[
		|C_{xx'}| = \sum_{\emptyset \neq W \subseteq \mathbf{t}} |H_W|.
		\]
		
		To compute the cardinality of \( H_W \), note that in every partition \( \pi \in H_W \), all vertices in the union \( \bigcup_{w \in W} p_w \) lie in the same block. Therefore, \( H_W \) is in bijection with the set of partitions of \( X_W \) into \( n - j \) blocks such that the associated graph contains no cycle passing through the collapsed vertex \( x_W \).
		
		Let \( n_W = |X_W| \) denote the number of vertices in the graph \( G_{\mathcal{R}}^W \). The corresponding partitions lie at level \( n_W - (n - j) \) in the lattice \( \Pi(X_W) \).
		
		To interrupt each of the \( t - |W| \) cycles in \( G_{\mathcal{R}}^W \) that pass through \( x_W \), we proceed as described above: construct a minimal cut \( U_{\gamma} \) by removing exactly one edge from each cycle.
		
		Thus, by the Inclusion-Exclusion principle, we obtain:
		\[
		|H_W| = \sum_u \sum_{\substack{\gamma \\ d(\gamma) = u(t - |W|)}} (-1)^{u+1} \bm{\pi}(X_W, n_W - (n - j), s - s_W, \mathcal{R} - U_{\gamma}).
		\]
		Let \( C_e \) be the set of partitions \( \pi \in \Pi(X, j, s, \mathcal{R})_{xx'} \) for which the \emph{only} path in \( G_{J_{\pi}} \) connecting \( x \) and \( x' \) is the edge \( e \).
		
		We claim that \( \pi \in C_e \) if and only if for every \( v \in \mathbf{t} \),
		\[
		|p_v - \mathcal{R}_{\pi}| \geq 2,
		\]
		where \( \mathcal{R}_{\pi} := \mathcal{A}(\pi) \cap \mathcal{R} \) denotes the set of atoms used in the decomposition of \( \pi \) that belong to \( \mathcal{R} \).
		
		Indeed, suppose \( |p_v - \mathcal{R}_{\pi}| = 0 \) for some \( v \in \mathbf{t} \). Then \( p_v \subseteq J_{\pi} \), and so \( G_{J_{\pi}} \) contains a path from \( x \) to \( x' \) different from \( e \), contradicting the assumption that \( \pi \in C_e \).
		
		Similarly, if \( |p_v - \mathcal{R}_{\pi}| = 1 \) for some \( v \in \mathbf{t} \), then all but one edge of \( p_v \) are present in \( J_{\pi} \), and a redundant edge could be added to complete a cycle passing through both \( x \) and \( x' \), again contradicting \( \pi \in C_e \).
		
		However, if \( |p_v- \mathcal{R}_{\pi}| \geq 2 \) for all \( v \in \mathbf{t} \), then even though redundant edges might exist, completing a cycle would change the block structure of \( \pi \), yielding a different partition. Thus, such \( \pi \) must lie in \( C_e \).
		
		To compute the cardinality of \( C_e \), we proceed by compressing the edge \( e \) and analyzing the resulting graph \( G_{\mathcal{R}}^{\{e\}} \), where \( e \) is replaced by a single vertex \( x_e \). In this graph, we must interrupt every cycle passing through \( x_e \) \emph{twice} in order to prevent alternate connections between \( x \) and \( x' \).
		
		Let \( U_{\gamma} \) be a minimal set constructed by selecting exactly two edges from each such cycle in \( G_{\mathcal{R}}^{\{e\}} \). Then,
		\[
		C_e = \bigcup_{\gamma} \Pi(X_e, j - 1, s - 1, \mathcal{R}' - U_{\gamma}),
		\]
		where \( X_e \) is the vertex set of the compressed graph and \( \mathcal{R}' \) denotes the set of atoms corresponding to the edges in \( G_{\mathcal{R}}^{\{e\}} \).
		
		Applying the Inclusion-Exclusion principle, we obtain:
		\[
		|C_e| = \sum_u \sum_{\substack{\gamma \\ d(\gamma) = u(t - 1)}} (-1)^{u+1} \bm{\pi}(X_e, j - 1, s - 1, \mathcal{R}' - U_{\gamma}).
		\]
		Therefore, since \( \Pi(X, j, s, \mathcal{R})_{xx'} = C_{xx'} \sqcup C_e \), the computations above yield the total cardinality of \( \Pi(X, j, s, \mathcal{R})_{xx'} \).
		
		Finally, to compute the cardinality of \( \Pi(X, j, s, \mathcal{R}) \), the set of partitions in which \( x \) and \( x' \) lie in \emph{different blocks}, we must interrupt all paths connecting \( x \) and \( x' \).
		
		Let \( U_{\alpha} \) be a minimal cut as defined earlier, obtained by removing exactly one edge from each such path. Then, using the Inclusion-Exclusion principle again, we obtain:
		\[
		|\Pi(X, j, s, \mathcal{R})| = \sum_u \sum_{\substack{\alpha \\ d(\alpha) = ut}} (-1)^{u+1} \bm{\pi}(X, j, s, \mathcal{R}' - U_{\alpha}).
		\]
		This completes the proof.
	\end{proof}

\end{document}
